\chapter{Dirac 作用量的 $O(a)$ 改进}
% Section 13 of R. Gupta, "Introduction to Lattice QCD" (hep-lat/9807028).

最早试图通过消除 $O(a)$ 修正来改进威尔逊费米子作用量的是 Hamber 和 Wu [102]。他们在威尔逊作用量中加入了一个双链项
\begin{align}
\mathcal{S}_{HW}
&= A_W + \frac{r\kappa}{4a} \sum_{x,\mu} \bar{\psi}^L(x)
   \left[ U_\mu (x) U_\mu (x+\hat{\mu}) \psi^L (x+2\hat{\mu})
   + U_\mu^\dagger (x-\hat{\mu}) U_\mu^\dagger (x-2\hat{\mu}) \psi^L (x-2\hat{\mu}) \right] \nonumber\\
&\equiv \sum_{x,y} \bar{\psi}_L (x) M_{xy}^{HW} \psi_L (y) .
\label{eq:HWaction}
\end{align}

容易通过泰勒展开检验，在经典层次上，$\mathcal{S}_W$ 中的 $O(a)$ 赝象被 HW 项所抵消。可以通过在微扰论中计算相对系数来计入量子修正。事实证明这一作用量并未得到深入研究——在 1980 年代是因为缺乏计算能力，而在 1990 年代则是因为 Sheikholeslami-Wohlert（简称 SW 或三叶草）作用量 [20] 的优点。

\section{Sheikholeslami-Wohlert（三叶草）作用量}

SW 提出在 $\mathcal{S}_W$ 中再加入一个维数为五的算符，即磁矩项，以消除 $O(a)$ 赝象而不破坏 Wilson 对加倍子的处理。所得 SW 作用量为
\begin{equation}
\mathcal{S}_{SW} \ = \ \mathcal{S}_W -
   \frac{i a C_{SW} \kappa r}{4}\ \bar{\psi}(x) \sigma_{\mu\nu} F_{\mu\nu} \psi(x) \ .
\label{eq:SWact}
\end{equation}

该作用量对在壳物理量是 $O(a)$ 改进的这一事实不能仅由泰勒展开推出。要证明改进，需要枚举所有维数为五的算符，利用运动方程消去其中一些，并通过重整化 $g^2 \to g^2 (1+b_g ma)$ 和 $m \to m (1+b_m ma)$ 将余下的两个（$m^2 \bar{\psi}\psi$ 与 $m F^{\mu\nu} F_{\mu\nu}$）吸收掉。SW 作用量的构造已由 L\"uscher 教授在本暑期学校中详细讲解，因此这里不再重复，而是请读者参阅他的讲义。我只简要总结用微扰论和非微扰方法调节 $C_{SW}$ 的结果。下面列出曾在模拟中使用过的 $C_{SW}$ 微扰改进值。

\begin{center}
\begin{tabular}{ll}
$C_{SW} = 1$                                & 树级改进 \\
$C_{SW} = 1 + 0.2659(1) g_0^2$              & 1 圈改进 [103] \\
$C_{SW} = 1/u_0^3$                          & 蝌蚪图改进 \\
$C_{SW} = 1/u_0^3 (1 + 0.0159 g_I^2)$       & 1 圈蝌蚪图改进 \\
\end{tabular}
\end{center}

在 1 圈 TI 结果中，$u_0$ 取为方格的四次方根。作为练习，请自行验证：使用朗道规范下链的期望值所得到的数值（$1.6$）更接近下面给出的非微扰改进结果（$g=1$ 时为 $1.77$）。对于上述每一个 $C_{SW}$ 值，误差都从 $O(a)$（威尔逊作用量）降低到至少 $O(\alpha_s a)$。

最近令人振奋的进展（L\"uscher 也作了评述）是：可以利用轴矢沃德恒等式 $C_{SW}$ 非微扰地调节 $C_{SW}$，从而去除全部 $O(a)$ 赝象，前提是同时改进耦合常数 $g^2$、夸克质量 $m_q$ 和流 [104]。对于采用威尔逊规范作用量的淬火理论，已在多个 $g$ 值下给出了 $C_{SW}$ 的估计 [104,105]。这些数据可用 [105] 拟合
\begin{equation}
C_{SW} = \frac{1 - 0.6084 g^2 - 0.2015 g^4 + 0.03075 g^6}{1 - 0.8743 g^2} \ .
\label{eq:CSWnonpert}
\end{equation}
在 $6/g^2 > 5.7$ 时成立。

SW 作用量的优点是它定域，且使微扰论仍可处理 [106]。加上三叶草项在威尔逊费米子模拟中只增加约 $\sim 15\%$ 的计算量。固定夸克质量（$M_\pi/M_\rho = 0.7$）时，使用 SW 作用量相对于威尔逊作用量所获改进的一个例子是：在 $a=0.17$ 费米处，$M_N/\sqrt{\sigma}$ 与 $M_{\rm vector}/\sqrt{\sigma}$ 偏离 $a=0$ 值的程度由 $30-40\%$ 降至 $\sim 3-5\%$ [105]。许多合作组正在对该改进做进一步的检验与评估。我预期在 LATTICE 98 会议上将报告若干淬火结果，这些结果应于 1998 年底前出现在电子预印本库中。

\section{D234c 作用量}

如前所述，离散化误差有两个来源：（i）用有限差分近似导数，以及（ii）量子修正，例如来自格点理论中不存在的动量 $> \pi/a$ 模的量子修正。Alford、Klassen 和 Lepage 提出了改进的规范与 Dirac 作用量，通过经典改进减小（i），通过蝌蚪图改进减小（ii）[107,79]。他们针对各向同性格点提出的最新方案称为 \textbf{D234c 作用量}，因为它涉及二阶、三阶和四阶导数，其形式为
\begin{align}
M_{D234c}
&= m(1+0.5 r a m) + \sum_\mu \left\{ \gamma_\mu \Delta_\mu^{(1)}
   - \frac{C_3}{6} a^2 \gamma_\mu \Delta_\mu^{(3)} \right\} \nonumber\\
&\quad + r \sum_\mu \left\{ - \frac{1}{2} a \Delta_\mu^{(2)}
   - \frac{C_F}{4} a \sum_\nu \sigma_{\mu\nu} F_{\mu\nu}
   + \frac{C_4}{24} a^2 \Delta_\mu^{(4)} \right\} \ .
\label{eq:D234caction}
\end{align}

这里 $\Delta_\mu^{(n)}$ 是 $n$ 阶格点协变导数（对称且延伸最短）。正比于 $C_3$ 的项恰是消除 $\Delta_\mu^{(1)}$ 中领头离散化误差所需（见式~7.6）。正比于 $r$ 的各项由场重定义产生，因此代表一个冗余算符。它们从而不改变朴素作用量对在壳物理量的 $O(a)$ 改进。其中三项分别为：威尔逊项（再次用来消除加倍子）、三叶草项（对 $F_{\mu\nu}$ 他们使用 $O(a^2)$ 改进版本，如 SW 作用量中用来消除 Wilson 处理所引入 $O(a)$ 误差那样）以及一个 $a^3$ 修正。为给经典值 $C_3 = C_F = C_4 = 1$ 计入量子修正，他们求助于蝌蚪图改进，即把所有出现在 $M_{D234c}$ 中的链都除以 $u_0$，其中 $u_0$ 取为朗道规范下的平均链。

量子修正通常也会诱导出额外的算符，不过在微扰论中这些算符从 $O(\alpha_s)$ 阶才开始出现。在蝌蚪图改进之后，这类算符的 $\alpha_s$ 系数预计都很小，因此人们希望其贡献可以忽略。其思想/动机类似于在式~12.14 给出的蝌蚪图改进的 Luscher-Weisz 纯规范作用量中略去小的 $\mathcal{L}_2^{(6)}$ 项。

对于其中两个系数 $C_F$ 和 $C_3$，作者对蝌蚪图改进的精度做了非微扰检验 [79]。由于只有算符 $\Delta_\mu^{(3)}$ 在 $O(a^2)$ 阶破坏转动对称性，因此可以通过研究色散关系的转动不变性来非微扰地调节 $C_3$。类似地，三叶草项和威尔逊项在 $O(a)$ 阶破坏手征对称性，因此它们的相对系数 $C_F$ 可以通过研究轴矢沃德恒等式来调节。他们的检验结果表明：即使在 $a = 0.4$ 费米处，对高达 $p = 1.5/a$ 的动量，谱与色散关系的误差也约为 $\approx 10\%$！作者正在把对 $D234$ 作用量的分析推广到各向异性格点，正如将在第 18.4 节讨论的，各向异性格点在胶球质量的计算中已经带来了丰硕的回报。

\section{NRQCD 作用量}

在重夸克模拟中，除了 $O(pa)$ 与 $O(\Lambda_{\rm QCD} a)$ 误差之外，还必须担心 $O(m_H a)$ 误差。对于粲夸克和底夸克，在可合理模拟的格距范围内这些误差为 $O(1)$。Lepage 及其合作者认识到，重-重或重-轻系统中的重夸克在很好的近似下是非相对论的，因而发展了一种对 $m_H a > 1$ 成立的非相对论表述（NRQCD）[108]。在这种方法中，两分量 Pauli 旋量的作用量由 Foldy-Wouthuysen 变换得到，相对论与离散化修正作为 $1/M$ 和 $p/M$ 的系统展开加入。NRQCD 作用量的一个版本为 [109]
\begin{align}
\mathcal{L}
&= \bar{\psi}_t \psi_t \nonumber\\
&\quad - \bar{\psi}_t
   \left( 1 - \frac{a \delta H}{2} \right)_t
   \left( 1 - \frac{a H_0}{2n} \right)^n_t
   U_4^\dagger
   \left( 1 - \frac{a H_0}{2n} \right)^n_{t-1}
   \left( 1 - \frac{a \delta H}{2} \right)_{t-1} \psi_{t-1},
\label{eq:nrqcdact}
\end{align}
其中 $\psi$ 是两分量 Pauli 旋量，$H_0$ 是非相对论动能算符，
\begin{equation}
H_0 = - \frac{\Delta^{(2)}}{2 M_0} \ .
\label{eq:nrqcdH0}
\end{equation}
$\delta H$ 包含相对论与有限格距修正，
\begin{align}
\delta H
&= - \frac{g c_1}{2 M_0} \boldsymbol{\sigma}\cdot\boldsymbol{B} \nonumber\\
&\quad + \frac{i g c_2}{8 (M_0)^2}
   \left( \boldsymbol{\Delta}\cdot\boldsymbol{E} - \boldsymbol{E}\cdot\boldsymbol{\Delta} \right)
   - \frac{g c_3}{8 (M_0)^2}
   \boldsymbol{\sigma}\cdot
   \left( \boldsymbol{\Delta}\times\boldsymbol{E} - \boldsymbol{E}\times\boldsymbol{\Delta} \right) \nonumber\\
&\quad - \frac{c_4 (\Delta^{(2)})^2}{8 (M_0)^3}
   + \frac{c_5 a^2 \Delta^{(4)}}{24 M_0}
   - \frac{c_6 a (\Delta^{(2)})^2}{16 n (M_0)^2} \ .
\label{eq:NRQCDdeltaH}
\end{align}
这里 $\boldsymbol{E}$ 和 $\boldsymbol{B}$ 是电场和磁场，$\boldsymbol{\Delta}$ 与 $\Delta^{(2)}$ 分别是空间规范协变格点导数和拉普拉斯算符，$\Delta^{(4)}$ 是连续算符 $\sum_i D_i^4$ 的格点版本。参数 $n$ 的引入是为了消除重夸克传播子中由最高动量模引起的不稳定性 [108]。同样，树级系数 $c_i = 1$ 可以通过把所有链矩阵按 $u_0$ 重标度而进行蝌蚪图改进。

这一表述已非常成功地用于从 $\upsilon$ 谱计算 $\alpha_s$ [110]（如第 19 节所讨论）、计算重-重与重-轻谱 [111]、以及重-轻衰变常数 $f_B$ 和 $f_{B_s}$ [109]。NRQCD 的缺点是它是一种有效理论，不能取 $a \to 0$ 极限。事实上，为模拟粲夸克和底夸克，必须在 $1 \lesssim 1/a \lesssim 2.5$ GeV 的狭窄窗口内工作，以保持 $m_H a > 1$。好消息是，当前的数据 [110,111,109] 表明，在式~13.7 所包含的各项下，离散化与相对论修正很小且处于可控范围。

\section{完美作用量}

完美 Dirac 作用量的发展尚未完成。其方法与规范作用量所用相同，见 [90] 中的讨论。瓶颈在于，鞍点积分给出的作用量即使在自由场极限下也涉及许多格点，且任意一对格点之间的规范连接是大量路径的平均 [112]。Degrand 及其合作者 [100] 已对这种截断到 $3^4$ 超立方体的作用量做了首批检验。从这些检验中尚不清楚，与其他离散化方案相比，这种改进是否值得为模拟这样的作用量付出额外代价。希望不久将找到更好的调节与检验此类作用量的方法。

\section{Fermilab 作用量}

Fermilab 小组提出了一种「Heavy Wilson」费米子作用量 [113]。该作用量意在介于轻夸克的 $O(a)$ 改进 SW-三叶草作用量与重夸克 NRQCD 作用量之间。其目标不仅是消除 $O(a)$ 离散化修正，还要计入全部幂次的 $ma$ 误差，从而使粲偶素和底偶素系统能以与轻-轻物理相同的可靠性进行模拟。该改进尤其针对改善重-轻矩阵元。到目前为止还没有用所提出作用量的检验报告，因此我不在此复述该作用量，而是请感兴趣的读者参阅 [113] 了解细节。

\section{小结}

与改进规范作用量的讨论相类似，既然上述每一种表述都在淬火格点数据中展现出了改进的切实证据，问题就在于：哪一作用量最适合做全面研究？在我看来，决定性因素将是动力学模拟，因为前置因子和标度指数（$L^{\sim 10}$）使得在没有改进的情况下，在当今计算机上更新背景组态的成本高得难以承受。使用改进作用量的动力学模拟才刚刚开始，因此还需要几年时间才能产生足够的数据，以判断不同方法的优缺点。

最后，我应该提及一种近来新颖的方法——畴壁费米子——用于改进威尔逊费米子的手征性质 [114]。在这种方法中，系统（畴）是五维的，费米子生活在构成相对边界（壁）的四维曲面上。两面壁上的四维规范场完全相同，而费米子零模具有相反的手征性。这些模之间的交叠被第五维方向上的分离指数地压低。因此，手征对称性破缺的贡献也受到抑制，通过相加两个畴壁的贡献，可以得到对 QCD 这样的类矢量理论具有良好手征性质的结果。弱矩阵元的首批模拟显示出良好的手征性质，总体上看起来非常有前景 [115]。由于本讲义未能充分涵盖这一令人振奋的方法，感兴趣的读者请参阅 [115] 及其中的参考文献。
